% ==============================================================================
% SECCION I: INTRODUCCION
% ==============================================================================
\section{Introducción}\label{sec:introduccion}

\subsection{Planteamiento y Relevancia del Problema}
El péndulo invertido rotacional de Furuta, introducido por Furuta \textit{et al.}~\cite{furuta1991} y sistematizado analíticamente por Cazzolato y Prime~\cite{cazzolato2011}, constituye un banco de pruebas fundamental en control digital e ingeniería mecatrónica. Su relevancia científica y formativa radica en que concentra, en una planta de laboratorio compacta, la transición entre maniobras no lineales de gran amplitud ---para elevar la barra desde el reposo inferior--- y la regulación fina alrededor de un punto de equilibrio inestable. Esta dualidad exige integrar modelado electromecánico riguroso, estimación discreta de velocidad libre de distorsión y ejecución determinista en tiempo real bajo restricciones estrictas de energía y carrera angular.

Aunque el problema de estabilización se resuelve con relativa holgura sobre bancos industriales equipados con motores de accionamiento directo, etapas de potencia MOSFET y encoders ópticos de alta cuenta, el panorama cambia sustancialmente al trasladar la planta hacia prototipos académicos reproducibles y de bajo costo. En plataformas manufacturadas con polímeros impresos y electrónica comercial accesible convergen de manera simultánea cuatro desafíos físicos y numéricos críticos:
\begin{enumerate}
    \item \textbf{Subactuación intrínseca SIMO:} una única entrada escalar de tensión $u(t)$ debe gobernar simultáneamente dos grados de libertad dinámicamente acoplados, a saber, el desplazamiento acimutal del brazo horizontal $\theta$ y la deflexión angular de la barra pendular $\alpha$.
    \item \textbf{Inestabilidad hiperbólica y fase no mínima:} en la vertical invertida ($\alpha = \pi$), la planta presenta un polo en el semiplano derecho que provoca divergencia rápida en lazo abierto, acompañado de dinámica de fase no mínima en la regulación de $\theta$; en consecuencia, para corregir una desviación del brazo sin perder la verticalidad, el actuador debe acelerar brevemente en sentido opuesto al error de posición.
    \item \textbf{No linealidades severas de hardware económico:} el empleo de un motor con reductor planetario $80:1$ introduce fricción seca mixta y holgura mecánica (\textit{backlash}), mientras que el puente~H bipolar L298N impone caída de tensión en su etapa Darlington y una zona muerta asimétrica acusada, factores que degradan la autoridad de control a baja señal y desencadenan ciclos límite sostenidos.
    \item \textbf{Ruido de cuantización en derivación discreta:} al estimar las velocidades angulares por diferenciación numérica a \qty{1.0}{\kilo\hertz} a partir de encoders magnéticos de 12~bits, los escalones de cuantización amplifican componentes de alta frecuencia capaces de excitar resonancias estructurales de la barra y del tren reductor si no se atenúan con mínimo retardo de fase.
\end{enumerate}

Ante la interacción de estas limitaciones, las estrategias lineales convencionales ajustadas sobre modelos ideales suelen oscilar o perder el equilibrio al enfrentarse al prototipo físico real. En este contexto, el presente estudio busca responder la siguiente pregunta de investigación: \textit{¿cómo diseñar, identificar experimentalmente e implementar un esquema de control digital y filtrado determinista a \qty{1.0}{\kilo\hertz} capaz de lograr levantamiento autónomo rápido y estabilización continua robusta sobre una plataforma impresa en PLA con actuador L298N y encoders magnéticos AS5600, venciendo sus no linealidades severas sin recurrir a sensores ópticos ni puentes MOSFET?}

\subsection{Objetivos y Estructura del Artículo}
El objetivo general de este trabajo es diseñar, caracterizar experimentalmente e implementar en tiempo real una arquitectura de control digital determinista a \qty{1.0}{\kilo\hertz}, apoyada en identificación empírica, filtrado selectivo adaptativo y supervisión mediante gemelo digital~3D, capaz de alcanzar el levantamiento autónomo rápido y la estabilización invertida continua de un péndulo de Furuta fabricado en PLA sobre hardware accesible.

Para dar cumplimiento a este propósito y responder a la pregunta planteada, se establecen cuatro objetivos específicos medibles:
\begin{enumerate}
    \item Diseñar y fabricar una plataforma mecatrónica modular en PLA bajo criterios de diseño para fabricación aditiva sin soportes (DfAM, $\text{SUP}=0$) con verificación booleana libre de interferencias, e instrumentarla en un microcontrolador STM32F446RE a \qty{1.0}{\kilo\hertz} con doble bus Soft-$\text{I}^2\text{C}$ y modulación PWM a \qty{20.0}{\kilo\hertz}.
    \item Caracterizar e identificar experimentalmente en banco real la zona muerta asimétrica del actuador, la función de transferencia motor--reductor y los parámetros pendulares acoplados.
    \item Sintetizar e implementar la arquitectura híbrida \texttt{M4 Ultra-Silk} ---integrada por levantamiento resonante en cuadratura, lazo dual paralelo ($\text{PD}_\alpha$ y $\text{PID}_\theta$), cascada de filtrado \textit{Dual-Notch} adaptativo, compensación continua de zona muerta y servo de retorno con reinicio de Clegg y freno activo--- acoplada en tiempo real a un gemelo digital~3D en MATLAB.
    \item Validar cuantitativamente el desempeño en lazo cerrado mediante comparación cruzada entre el modelo CAD/teórico, la simulación no lineal a \qty{1.0}{\kilo\hertz} y los ensayos físicos en tiempo real, evaluando tiempos de levantamiento, error vertical RMS, esfuerzo de actuación y recuperación ante perturbaciones.
\end{enumerate}

El resto del artículo se organiza de la siguiente manera: la Sección~\ref{sec:marco_teorico} establece el marco teórico, la revisión de antecedentes y el modelado dinámico matricial del sistema; la Sección~\ref{sec:metodologia} detalla la metodología de diseño mecánico, instrumentación electrónica, identificación experimental, síntesis de control, integración del gemelo digital~3D y desarrollo del simulador cinemático interactivo~3D; la Sección~\ref{sec:resultados} presenta los resultados de simulación y validación experimental en planta real; la Sección~\ref{sec:discusiones} interpreta críticamente los hallazgos frente a los objetivos y la literatura; finalmente, la Sección~\ref{sec:conclusiones} expone las conclusiones y el trabajo futuro.

Con el fin de garantizar la trazabilidad y reproducibilidad integral de esta investigación, la totalidad de los artefactos del proyecto ---modelos CAD paramétricos en FreeCAD~1.1, placas de impresión~3D sin soportes, firmware embebido en C para STM32, conjuntos de datos experimentales CSV y guiones de identificación, el gemelo digital~3D en MATLAB, el simulador interactivo~3D de cinemática matricial (\texttt{simulacion cinematica/\allowbreak index.html}) y los archivos fuente LaTeX del artículo--- se encuentra publicada en acceso abierto en el repositorio oficial de GitHub: \url{https://github.com/Nocudo/Pendulo-Furuta-Control-Digital}.
